% =========================================================
% APÊNDICES
% FINALIDADE: reunir material ELABORADO PELO(A) AUTOR(A) que complementa
%   o texto (questionários, roteiros, termos, quadros extensos).
%   Material de outra autoria vai em anexos.tex.
% ESCREVA: cada apêndice começa com \apendice{TÍTULO EM MAIÚSCULAS}.
%   A letra (A, B, C...) e a entrada no sumário são automáticas.
% OPCIONAL: sem apêndices, apague este arquivo; o elemento é omitido.
% CUIDADO: cite cada apêndice no texto, por exemplo:
%   conforme o Apêndice~\ref{apd:instrumento}.
% EXEMPLO de segundo apêndice (retire os % para usar):
%   \apendice{ROTEIRO DE ENTREVISTA}
%   \label{apd:roteiro}
%   Texto do apêndice.
% =========================================================

\apendice{INSTRUMENTO DE COLETA DE DADOS}
\label{apd:instrumento}

Elemento opcional. Apresente, em cada apêndice, materiais elaborados pelo(a) autor(a), como questionários, roteiros de entrevista, protocolos de observação, termos de consentimento e quadros complementares.

% Para inserir um documento em PDF, salve-o na pasta 5-figuras/ e
% retire o % da linha abaixo, ajustando o nome do arquivo.
% pages=- insere todas as páginas; pages={1-3} insere só as páginas 1 a 3.
% \includepdf[pages=-,pagecommand={}]{5-figuras/nome-do-arquivo.pdf}
